% TOP_PROBLEM: {"id": 1183, "title": "Artin's Conjecture on Primitive Roots", "category_id": 1, "category": {"id": 1, "name": "number_theory", "display_name": "Number Theory", "description": "Properties of integers, prime numbers, Diophantine equations.", "slug": "number-theory", "order_index": 1, "created_at": "2026-07-31T15:26:25.670Z"}, "status": "open"}
% FIRST_POSTED: 2026-09-27
% !TeX program = lualatex
% Compile twice: lualatex open_problems_500.tex
% All references and prose are embedded; no BibTeX database or external inputs.
\documentclass[11pt,a4paper]{article}
\usepackage[margin=24mm,headheight=14pt]{geometry}
\usepackage{fontspec}
\setmainfont{Latin Modern Roman}
\setsansfont{Latin Modern Sans}
\setmonofont{Latin Modern Mono}
\usepackage{amsmath,amssymb,mathtools}
\usepackage{unicode-math}
\setmathfont{Latin Modern Math}
\usepackage{microtype}
\usepackage{enumitem,xurl,needspace}
\usepackage[dvipsnames]{xcolor}
\usepackage{hyperref}
\hypersetup{unicode=true,colorlinks=true,linkcolor=MidnightBlue,urlcolor=MidnightBlue,
 pdftitle={500 Mathematical Problem Statements},pdfauthor={Source-annotated compilation},bookmarksnumbered=true}
\usepackage{bookmark}
\usepackage{tocloft}
\setlength{\cftsecnumwidth}{3em}
\cftsetpnumwidth{2.5em}
\cftsetrmarg{4em}
\usepackage{titlesec}
\titleformat{\section}{\Large\bfseries\raggedright}{\thesection}{0.7em}{}
\usepackage{fancyhdr}
\pagestyle{fancy}\fancyhf{}
\fancyhead[L]{\small\sffamily Mathematical problem statements}
\fancyhead[R]{\small Source edition 22}
\fancyfoot[C]{\thepage}
\setlength{\parindent}{0pt}
\setlength{\parskip}{5pt plus 1pt}
\setlength{\emergencystretch}{3em}
\setcounter{tocdepth}{1}
\setcounter{secnumdepth}{1}
\setlist[itemize]{leftmargin=3em,itemsep=5pt,topsep=4pt,parsep=0pt}
\allowdisplaybreaks
\newcommand{\N}{\mathbb{N}}
\newcommand{\Z}{\mathbb{Z}}
\newcommand{\Q}{\mathbb{Q}}
\newcommand{\R}{\mathbb{R}}
\newcommand{\C}{\mathbb{C}}
\newcommand{\F}{\mathbb{F}}
\newcommand{\A}{\mathbb{A}}
\newcommand{\PP}{\mathbb P}
\newcommand{\E}{\mathbb{E}}
\newcommand{\eps}{\varepsilon}
\newcommand{\dd}{\,\mathrm d}
\newcommand{\sref}[2]{\hyperlink{source:#1:#2}{[S#2]}}
\newcommand{\eref}[2]{\hyperlink{extra:#1:#2}{[E#2]}}
% Turn the next switch off for a shorter reading copy. The quotations preserve
% every exactTarget field, including qualifications omitted from a short summary.
\newif\ifshowcatalogue
\showcataloguetrue
\newcommand{\cataloguescope}[1]{\ifshowcatalogue
 \paragraph{Exact scope in the supplied catalogue.}
 \begin{quote}\small #1\end{quote}\fi}

% Compile twice with: lualatex 1183.tex
\numberwithin{equation}{section}
\hypersetup{pdftitle={Research attempt -- problem 1183},pdfauthor={Research notebook; original compilation retained}}
\fancyhead[L]{\small\sffamily Research notebook}
\fancyhead[R]{\small\sffamily Problem 1183 | 27 September 2026}
\begin{document}
\setcounter{section}{0}
% ================================================================
% problemId: 1183
% Canonical problem ID: 1183
% exactTarget SHA-256: 5d51abfe129974858a8948e1ae8bd1a9c674175a4eb964e45336020ff729a07c
\clearpage
\section*{\texorpdfstring{Artin's Conjecture on Primitive Roots}{Artin's Conjecture on Primitive Roots}}\label{1183}
\subsection{Definitions and mathematical statement}
For a prime $p\nmid a$, let $\operatorname{ord}_p(a)$ be the least $m>0$ with $a^m\equiv1\pmod p$. Then $a$ is a primitive root modulo $p$ exactly when $\operatorname{ord}_p(a)=p-1$. For each integer $a\ne-1$ which is not a square, the target requires a positive prime-relative density
\[
 \lim_{x\to\infty}\frac{\#\{p\le x:\operatorname{ord}_p(a)=p-1\}}{\pi(x)}=\delta(a)>0.
\]
The Artin density is expressed in the usual inclusion-exclusion convention by
$\delta(a)=\sum_{m\ge1}\mu(m)/[\Q(\zeta_m,a^{1/m}):\Q]$, equivalently its corrected Artin-product form. The algebraic degrees encode dependencies between prime-power obstructions. The density is among primes, not among all positive integers.
\par\smallskip\textit{Formulation and concept references:} \sref{1183}{1}.
\cataloguescope{Prove or disprove that every integer a other than -1 that is not a perfect square is a primitive root modulo p for a positive natural density of primes p, with the density given by Artin's predicted correction factor.}
\subsection{Short English statement}
Does every admissible integer generate the multiplicative group modulo a positive proportion of primes, with precisely Artin's predicted density?
% BEGIN RESEARCH ADDITION ATTEMPT-197
\subsection{Research attempt}

\paragraph{Current frontier.}
Hooley proved the predicted density conditional on suitable generalized Riemann hypotheses for the associated Kummer fields. Unconditional results for sets of possible primitive roots or for averaged variants do not establish the prescribed density for every fixed admissible $a$ \eref{1183}{1}. The catalogue's 2024 source concerns a two-variable variant; its introductory formulation is compatible with, but does not solve, the fixed-integer target \sref{1183}{1}. We pursue finite inclusion--exclusion plus a uniform tail estimate, keeping the entanglement correction in the actual field degrees.

\paragraph{Attack route.}
Fixed-field Chebotarev gives every finite collection of power-residue tests. The missing interchange is between infinitely many such tests and the prime-counting limit. One route assumes effective Chebotarev with GRH-quality errors; another seeks cancellation after averaging over the Kummer degree. We isolate the needed tail and prove an elementary bound for its large-prime-divisor portion. This leaves a specific middle range, not an unspecified appeal to ``independence of residue conditions.''

\paragraph{Attempt.}
For $p\nmid a$, define the residual index
\[
 i_p(a)=\frac{p-1}{\operatorname{ord}_p(a)}.
\]
For a prime $q$, $q\mid i_p(a)$ exactly when $q\mid p-1$ and $a^{(p-1)/q}\equiv1\pmod p$. Away from ramification this is equivalent to complete splitting of $p$ in
$L_q=\Q(\zeta_q,a^{1/q})$. For squarefree $m$, simultaneous splitting for its prime factors is splitting in $L_m=\Q(\zeta_m,a^{1/m})$: these fields form the corresponding compositum. The equivalences and degree corrections are standard Kummer ingredients in \eref{1183}{1}.

Let $P(y)=\prod_{q\le y}q$, and let $N_y(x)$ count primes $p\le x$, $p\nmid a$, with $\gcd(i_p(a),P(y))=1$. Finite inclusion--exclusion gives
\begin{equation}\label{eq197:finite}
 N_y(x)=\sum_{m\mid P(y)}\mu(m)\pi_{\rm split}(L_m,x)+O_{a,y}(1).
\end{equation}
All sums are finite for fixed $y$. Chebotarev therefore gives
\[
 \lim_{x\to\infty}\frac{N_y(x)}{\pi(x)}
 =\delta_y(a):=\sum_{m\mid P(y)}\frac{\mu(m)}{[L_m:\Q]}.
\]
The degree computation underlying Artin's corrected density yields
$\delta_y(a)\to\delta(a)>0$ for the admissible integers in the question; we retain this established algebraic computation, including its entanglement correction, as recorded in \eref{1183}{1}. A product assuming all fields linearly disjoint would not in general be the required constant.

Let $N_a(x)$ count actual primitive-root primes and set
\[
 B_a(x,y)=\#\{p\le x:p\nmid a,\ \exists q>y\text{ prime with }q\mid i_p(a)\}.
\]
Then $0\le N_y(x)-N_a(x)\le B_a(x,y)$. Consequently the sufficient missing estimate is
\begin{equation}\label{eq197:tail}
 \lim_{y\to\infty}\limsup_{x\to\infty}
       \frac{B_a(x,y)}{\pi(x)}=0.
\end{equation}
To verify sufficiency rigorously, fix $\varepsilon>0$, choose $y$ with $|\delta_y-\delta|<\varepsilon$ and limiting tail at most $\varepsilon$, and then take upper and lower limits in the preceding inequalities. Both limits for $N_a(x)/\pi(x)$ lie within $2\varepsilon$ of $\delta$. Sending $\varepsilon\to0$ proves the desired limit. This does not assert that \eqref{eq197:tail} follows from the already known fixed-$y$ limits.

\textbf{Large-$q$ bound.} Fix $B>0$. The number of primes $p\le x$ with a prime divisor
$q>\sqrt x(\log x)^B$ of $i_p(a)$ is $o(\pi(x))$.
The hypotheses on $a$ imply $|a|\ge2$. Such a prime has
\[
 \operatorname{ord}_p(a)\le\frac{p-1}{q}<M:=\frac{\sqrt x}{(\log x)^B}.
\]
Discarding the $O(\sqrt x)$ primes $p\le\sqrt x$, all remaining ones divide the nonzero integer
\[
 R_M=\prod_{1\le j\le\lfloor M\rfloor}|a^j-1|.
\]
Since $\log|a^j-1|\le j\log|a|+\log2$,
$\log R_M\ll_a M^2$. Each distinct prime above $\sqrt x$ contributes at least $\tfrac12\log x$, regardless of repeated divisibility among factors. Their number is at most
\[
 \frac{2\log R_M}{\log x}
 \ll_a\frac{x}{(\log x)^{2B+1}}=o\!\left(\frac{x}{\log x}\right).
\]
Together with the prime number theorem this proves the claim. Negative $a$ causes no problem because absolute values were used and no factor is zero.

Thus only the middle range
\[
 y<q\le\sqrt x(\log x)^B
\]
remains in \eqref{eq197:tail}. Controlling it uniformly would finish the density argument. Ignoring the power-residue condition and merely counting $p\equiv1\pmod q$ loses the extra suppression needed in this range. Applying a fixed-field asymptotic separately to a number of fields growing with $x$ supplies no uniform error bound. Hooley's conditional method addresses precisely such uniform issues; citing it without its analytic hypotheses would not give an unconditional proof \eref{1183}{1}.

\paragraph{Stress test.}
The finite inclusion--exclusion identity ignores only finitely many ramified primes for fixed $y$, not an uncontrolled set growing in the final limit. The large-$q$ product counts distinct prime divisors using logarithmic size, so it never assumes the factors $a^j-1$ are coprime. The tail bound is a sufficient statement; no unsupported assertion of equivalence is needed. The elementary argument does not establish the middle-range saving, and the exact small-prime order checks in the script do not justify an exchange of $x$ and $y$ limits. A proof of infinitude without the limiting constant would still be weaker than the printed target.

\paragraph{Outcome.}
\textbf{Outcome: Conditional reduction.}
The exact corrected density follows from a stated uniform tail estimate. Its very large prime-divisor range is bounded unconditionally by an explicit elementary product argument; the middle range remains unresolved. The derivation is a careful reconstruction of a standard strategy, not a new unconditional proof of Artin's conjecture or a claim that the correction factor can be omitted.

% END RESEARCH ADDITION ATTEMPT-197
\subsection{Sources}
\begin{itemize}\raggedright
\item[\textnormal{[S1]}]\hypertarget{source:1183:1}{} Journal of Number Theory 262 (2024), 161-185.
\url{https://www.sciencedirect.com/science/article/pii/S0022314X24000829}.
{\small\textit{Catalogue formulation source.}}
% BEGIN RESEARCH ADDITION REFERENCES-197
\item[\textnormal{[E1]}]\hypertarget{extra:1183:1}{} Pieter Moree, \emph{Artin\textquotesingle s primitive root conjecture---a survey}, Integers 12 (2012), 1305--1416; arXiv:math/0412262v2. Discussion of Hooley\textquotesingle s theorem, Kummer fields, corrected densities and the large-index sieve ranges.
\url{https://arxiv.org/html/math/0412262v2}.
{\small\textit{Research reference; consulted 27 September 2026.}}
% END RESEARCH ADDITION REFERENCES-197
\end{itemize}

\end{document}
